\documentclass[10pt,a4paper]{amsart}
\usepackage{amsmath,amssymb,amsthm,graphicx}
\usepackage{verbatim}
\newenvironment{reference}{\comment}{\endcomment}
%\newenvironment{reference}{}{}
\newenvironment{slogan}{\comment}{\endcomment}
\newenvironment{history}{\comment}{\endcomment}

% For commutative diagrams we use Xy-pic
\usepackage[all]{xy}

% We use 2cell for 2-commutative diagrams.
\xyoption{2cell}
% Theorem environments.
%
\theoremstyle{plain}
\newtheorem{theorem}[subsection]{Theorem}
\newtheorem{proposition}[subsection]{Proposition}
\newtheorem{lemma}[subsection]{Lemma}

\theoremstyle{definition}
\newtheorem{definition}[subsection]{Definition}
\newtheorem{example}[subsection]{Example}
\newtheorem{exercise}[subsection]{Exercise}
\newtheorem{situation}[subsection]{Situation}

\theoremstyle{remark}
\newtheorem{remark}[subsection]{Remark}
\newtheorem{remarks}[subsection]{Remarks}

\numberwithin{equation}{subsection}

% Macros
%
\def\lim{\mathop{\mathrm{lim}}\nolimits}
\def\colim{\mathop{\mathrm{colim}}\nolimits}
\def\Spec{\mathop{\mathrm{Spec}}}
\def\Hom{\mathop{\mathrm{Hom}}\nolimits}
\def\Ext{\mathop{\mathrm{Ext}}\nolimits}
\def\SheafHom{\mathop{\mathcal{H}\!\mathit{om}}\nolimits}
\def\SheafExt{\mathop{\mathcal{E}\!\mathit{xt}}\nolimits}
\def\Sch{\mathit{Sch}}
\def\Mor{\mathop{\mathrm{Mor}}\nolimits}
\def\Ob{\mathop{\mathrm{Ob}}\nolimits}
\def\Sh{\mathop{\mathit{Sh}}\nolimits}
\def\NL{\mathop{N\!L}\nolimits}
\def\CH{\mathop{\mathrm{CH}}\nolimits}
\def\proetale{{pro\text{-}\acute{e}tale}}
\def\etale{{\acute{e}tale}}
\def\QCoh{\mathit{QCoh}}
\def\Ker{\mathop{\mathrm{Ker}}}
\def\Im{\mathop{\mathrm{Im}}}
\def\Coker{\mathop{\mathrm{Coker}}}
\def\Coim{\mathop{\mathrm{Coim}}}

% Boxtimes
%
\DeclareMathSymbol{\boxtimes}{\mathbin}{AMSa}{"02}

%
% Macros for moduli stacks/spaces
%
\def\QCohstack{\mathcal{QC}\!\mathit{oh}}
\def\Cohstack{\mathcal{C}\!\mathit{oh}}
\def\Spacesstack{\mathcal{S}\!\mathit{paces}}
\def\Quotfunctor{\mathrm{Quot}}
\def\Hilbfunctor{\mathrm{Hilb}}
\def\Curvesstack{\mathcal{C}\!\mathit{urves}}
\def\Polarizedstack{\mathcal{P}\!\mathit{olarized}}
\def\Complexesstack{\mathcal{C}\!\mathit{omplexes}}
% \Pic is the operator that assigns to X its picard group, usage \Pic(X)
% \Picardstack_{X/B} denotes the Picard stack of X over B
% \Picardfunctor_{X/B} denotes the Picard functor of X over B
\def\Pic{\mathop{\mathrm{Pic}}\nolimits}
\def\Picardstack{\mathcal{P}\!\mathit{ic}}
\def\Picardfunctor{\mathrm{Pic}}
\def\Deformationcategory{\mathcal{D}\!\mathit{ef}}
\usepackage{cleveref}
\crefname{originallemma}{original lemma}{original lemmas}
\Crefname{originallemma}{Original lemma}{Original lemmas}
\makeatletter
\newlabel{lemma-relocalize}{{03EH}{source}{}{sourcepointer-1}{}}
\newlabel{lemma-relocalize@cref}{{[originallemma][0][]03EH}{[1][0][]source}{}{}{}}
\makeatother
\crefname{originalsection}{original section}{original sections}
\Crefname{originalsection}{Original section}{Original sections}
\makeatletter
\newlabel{section-sheafification}{{00W1}{source}{}{sourcepointer-2}{}}
\newlabel{section-sheafification@cref}{{[originalsection][0][]00W1}{[1][0][]source}{}{}{}}
\makeatother
\title{Maps of sheaves glue}
\author{The Stacks Project Authors}
\date{stacks-project:sites.tex:local-snapshot-20261012; GNU FDL 1.2 or later}
\begin{document}
\maketitle
\noindent\textit{Editorial note.} All delivered mathematical text is the author's own, copied from the pinned Stacks Project chapter file sites-cohomology.tex; no statement or proof step was rewritten, reordered or removed.\par\medskip
\noindent\textit{Editorial note.} The proof is the author's complete construction of the glued map; the sentence 'We omit the verification that these maps are compatible with the restriction maps' concerns the routine verification that the construction is functorial, and it is preserved verbatim inside the counted span.\par\medskip
\noindent\textit{Editorial note.} Every labelled reference occurring inside the counted spans is carried in assembly.external\_references as a non-load-bearing navigation pointer (same-chapter ones at their definition site, cross-chapter ones at the invocation inside the counted span), so no load-bearing mathematics is replaced by a stub. No bibliography citation occurs inside the counted spans.\par\medskip
\noindent\textit{Editorial note.} The checked chapter list for same-name collisions is the set of stacks chapters staged in this repository: chow, cohomology, divisors, etale-cohomology, flat, more-algebra, more-morphisms, morphisms, spaces-more-morphisms, varieties together with this card's chapter sites-cohomology.\par\medskip
\noindent\textit{Editorial note.} The preamble is the chapter's own retained minimal preamble (sources/stacks-sites-cohomology/retained-preamble.tex), a verbatim excerpt of the native preamble: the Xy-pic setup and the native theorem declarations with the mathematical macros.\par\medskip
\noindent\textit{Allowed standard prerequisites (editorial).}\begin{itemize}
\item Restriction of sheaves along a relocalization. Same chapter, Lemma lemma-relocalize. Supplies the restriction functors appearing in the statement.
\item The value of a sheaf on a covering. Same chapter, Section section-sheafification, where the equalizer describing the value on a covering is introduced. Supplies the sheaf condition used to define the glued map.
\end{itemize}\medskip
\ifdefined\hypertarget\hypertarget{sourcepointer-1}{}\fi
\noindent\textit{Source reference.} Original lemma 03EH: sources/stacks-sites/pinned-original-source.tex, source lines 5833--5833. Navigation pointer invoked at proof line 5833; the counted span contains the pointer only and the cited result is not reproduced inside it, so the pointer is registered as a non-load-bearing navigation stub in the established form for Stacks material. No bibliography citation is involved.\par\medskip
\ifdefined\hypertarget\hypertarget{sourcepointer-2}{}\fi
\noindent\textit{Source reference.} Original section 00W1: sources/stacks-etale-cohomology/pinned-original-source.tex, source lines 725--725. Navigation pointer invoked at proof line 5913; the counted span contains the pointer only and the cited result is not reproduced inside it, so the pointer is registered as a non-load-bearing navigation stub in the established form for Stacks material. No bibliography citation is involved.\par\medskip
\begin{lemma}
\label{lemma-glue-maps}
\begin{slogan}
Maps of sheaves glue.
\end{slogan}
Let $\mathcal{C}$ be a site.
Let $\{U_i \to U\}$ be a covering of $\mathcal{C}$.
Let $\mathcal{F}$, $\mathcal{G}$ be sheaves on $\mathcal{C}$.
Given a collection
$$
\varphi_i :
\mathcal{F}|_{\mathcal{C}/U_i}
\longrightarrow
\mathcal{G}|_{\mathcal{C}/U_i}
$$
of maps of sheaves such that for all $i, j \in I$ the maps
$\varphi_i, \varphi_j$ restrict to the same map
$\varphi_{ij} : \mathcal{F}|_{\mathcal{C}/U_i \times_U U_j} \to
\mathcal{G}|_{\mathcal{C}/U_i \times_U U_j}$
then there exists a unique map of sheaves
$$
\varphi :
\mathcal{F}|_{\mathcal{C}/U}
\longrightarrow
\mathcal{G}|_{\mathcal{C}/U}
$$
whose restriction to each $\mathcal{C}/U_i$ agrees with $\varphi_i$.
\end{lemma}

\begin{proof}
The restrictions used in the lemma are those of
Lemma \ref{lemma-relocalize}.
Let $V/U$ be an object of $\mathcal{C}/U$.
Set $V_i = U_i \times_U V$ and denote $\mathcal{V} = \{V_i \to V\}$.
Observe that $(U_i \times_U U_j) \times_U V = V_i \times_V V_j$.
Then we have
$\mathcal{F}|_{\mathcal{C}/U_i}(V_i/U_i) = \mathcal{F}(V_i)$
and
$\mathcal{F}|_{\mathcal{C}/U_i \times_U U_j}(V_i \times_V V_j/U_i \times_U U_j)
= \mathcal{F}(V_i \times_V V_j)$
and similarly for $\mathcal{G}$.
Thus we can define $\varphi$ on sections over $V$ 
as the dotted arrows in the diagram
$$
\xymatrix{
\mathcal{F}(V) \ar@{=}[r] &
H^0(\mathcal{V}, \mathcal{F}) \ar@{..>}[d] \ar[r] &
\prod \mathcal{F}(V_i)
\ar[d]_{\prod \varphi_i}
\ar@<1ex>[r] \ar@<-1ex>[r] &
\prod \mathcal{F}(V_i \times_V V_j) \ar[d]_{\prod \varphi_{ij}} \\
\mathcal{G}(V) \ar@{=}[r] &
H^0(\mathcal{V}, \mathcal{G}) \ar[r] &
\prod \mathcal{G}(V_i)
\ar@<1ex>[r] \ar@<-1ex>[r] &
\prod \mathcal{G}(V_i \times_V V_j)
}
$$
The equality signs come from the sheaf condition; see
Section \ref{section-sheafification} for the notation
$H^0(\mathcal{V}, -)$.
We omit the verification that these maps are compatible
with the restriction maps.
\end{proof}

\end{document}
